\subsection{Annexes}
\label{sec:Annexes}

\begin{frame}[c]{Annexes}
    \begin{itemize}
        \item Un résultat brut sans documentation (méthode, approche, hypothèses, valeurs d’entrée) doit passer directement à la poubelle. En complétant ce fichier, vous aurez déjà fait 90\% du travail.
        \item Un rendu suit les mêmes règles que l'archivage: seules les géométries d'entrée, le modèle de calcul, les outils de validation et de calcul doivent être conservés, sans fichiers de résultats, il sera toujours possible de les régénérer si nécessaire.
        \item En principe, il faut donc fournir un fichier .zip contenant ce document comme documentation, les fichiers nécessaires au solveur, ainsi que tous les scripts et/ou sources requis. Évidemment, les fichiers doivent être bien organisés.
        \item Sur cette diapositive, il est nécessaire de lister ces annexes et de spécifier l'auteur dans le cas où il s'agirait de quelqu'un d'autre. 
        \item Il n'y a aucun problème à utiliser le travail de quelqu'un d'autre, le problème serait de prétendre que nous en sommes l'auteur.
    \end{itemize}
\end{frame}